\section{Undecidable local reachability in universal QCA: a reduction}
\label{app:reachability_undecidability}

This appendix provides a concrete halting-to-reachability reduction establishing Theorem~\ref{thm:undecidable_reachability}.
Once a fixed local, reversible dynamics can simulate a universal machine and a halting event is recorded in a localized flag, the corresponding hitting predicate is undecidable by a direct reduction from halting.
For classical CA undecidability results in closely related settings see, e.g., \cite{Kari1990,Kari1994}.
For undecidability results about reachability properties of quantum-system models (quantum automata), see \cite{LiYing2014}.

\subsection{Formal predicate}
Fix a lattice $\ZZ$ and a local Hilbert space $\CC^d$ per site.
Let $U$ be a QCA update on $\bigotimes_{n\in\ZZ}\CC^d$.
For a finite region $K$ and a local projector $\Pi_K$ supported on $K$, define the \emph{local reachability predicate}
\begin{equation}
\mathsf{Reach}(U,\Pi_K;\psi)\ :=\ \Bigl[\exists t\ge 0:\ \langle\psi|U^{\dagger t}\Pi_K U^t|\psi\rangle>0\Bigr].
\end{equation}
Theorem~\ref{thm:undecidable_reachability} claims the existence of a fixed $(U,\Pi_K)$ for which $\mathsf{Reach}(U,\Pi_K;\psi_{M,x})$ is undecidable as a function of the encoded instance $(M,x)$.

\subsection{Step 1: reversible halting as a well-posed predicate}
The classical halting problem is undecidable \cite{Turing1936}.
To embed it into a unitary dynamics, it is convenient to work with a reversible (injective) machine model.
Reversible simulation of Turing machines is standard in reversible computation \cite{Bennett1973,Bennett1989}.
In particular, there exist universal reversible Turing machines whose halting behavior encodes the usual halting problem.

\subsection{Step 2: a fixed local, reversible dynamics simulating a universal reversible machine}
Let $M_\star$ be a fixed universal reversible Turing machine (URTM).
Fix a finite tape alphabet $\Gamma$, a finite internal state set $Q$, and a distinguished halting state $q_{\mathrm{halt}}\in Q$.
Encode the URTM configuration into a 1D lattice by using a finite local alphabet that stores:
(i) a tape symbol in $\Gamma$,
(ii) a head marker plus internal state in $Q$ at exactly one site, and
(iii) auxiliary registers used to make halting-event recording reversible \cite{Bennett1973,Bennett1989}.

\noindent One can implement one URTM step by a translation-invariant, finite-radius reversible local rule (equivalently, a reversible cellular automaton on the computational basis).
Promoting this reversible basis permutation to a unitary yields a QCA update $U$ which acts as the classical reversible dynamics on basis states.
Universal QCA constructions and QCA structure theorems are discussed in \cite{Watrous1995,SchumacherWerner2004,ArrighiGrattage2010,ArrighiNesmeWerner2011}.

\subsection{Step 3: a localized halting flag projector}
Given an instance $(M,x)$, construct an initial basis state $|\psi_{M,x}\rangle$ that encodes $M$ and input $x$ in a finite ``program/data'' region, with blank padding elsewhere.
Reserve a dedicated finite region $K$ containing a ``flag cell'' whose local basis includes states $|0\rangle$ and $|1\rangle$.

Design the simulation so that:
(i) if $M$ halts on $x$, then at the moment of halting the flag cell is set to $|1\rangle$ and remains $|1\rangle$ thereafter (this can be achieved without violating reversibility by recording the halting event in an ancilla and continuing with a reversible idle loop);
(ii) if $M$ does not halt, the flag cell remains $|0\rangle$ forever.

Let $\Pi_K$ be the local projector onto the subspace where the flag cell is in state $|1\rangle$ (tensored with identity on the rest of $K$).
Then
\begin{equation}
\exists t\ge 0:\ \langle\psi_{M,x}|U^{\dagger t}\Pi_K U^t|\psi_{M,x}\rangle>0
\end{equation}
holds if and only if $M$ halts on $x$.

\begin{proof}
Because the evolution is classical on the computational basis, for each $t$ the state $U^t|\psi_{M,x}\rangle$ is again a basis state.
Hence the quantity $\langle\psi_{M,x}|U^{\dagger t}\Pi_K U^t|\psi_{M,x}\rangle$ is in $\{0,1\}$ and equals $1$ iff the flag cell is in state $|1\rangle$ at time $t$.
By construction, the flag is triggered at some time iff the simulated computation halts.
\end{proof}

\subsection{Conclusion: undecidability}
If there were an algorithm deciding the reachability predicate above for all $(M,x)$, it would decide halting.
Therefore the local reachability predicate is undecidable.

\begin{remark}[Why this is compatible with unitarity]
The construction never modifies $U$ nonunitarily.
Undecidability arises because the predicate quantifies over unbounded time and because universal dynamics can encode arbitrarily long computations.
\end{remark}

\subsection{Robustness: a thresholded reachability predicate}
\label{app:reachability_robustness}
The bare predicate $\exists t:\langle\Pi_K\rangle>0$ is intentionally sharp, but it is not robust under arbitrary perturbations of the initial encoding: an arbitrarily small component in the flag subspace would make the inequality true.
One can, however, state a robust version that remains equivalent for perturbations bounded in trace distance.

\begin{lemma}[Projector expectation is Lipschitz in trace distance]
\label{lem:proj_lipschitz_trace}
For any two density operators $\rho,\sigma$ and any projector $\Pi$,
\begin{equation}
\bigl|\Tr(\Pi\rho)-\Tr(\Pi\sigma)\bigr|\ \le\ D(\rho,\sigma),
\qquad D(\rho,\sigma):=\tfrac12\|\rho-\sigma\|_1.
\end{equation}
\end{lemma}

\begin{proof}
This is standard: for any effect $0\le E\le \mathds{1}$ one has
$|\Tr(E(\rho-\sigma))|\le \tfrac12\|\rho-\sigma\|_1$; see, e.g., \cite{Watrous2018}.
\end{proof}

\begin{proposition}[Robust halting-to-reachability with a probability threshold]
\label{prop:reachability_threshold_robust}
Let $U$ and $\Pi_K$ be as in the reduction above and let $\rho_{M,x}:=|\psi_{M,x}\rangle\langle\psi_{M,x}|$ be the ideal basis-state encoding.
Assume the reduction is \emph{deterministic on the computational basis} so that for all $t$,
\begin{equation}
\Tr\!\bigl(\Pi_K\,U^t\rho_{M,x}U^{\dagger t}\bigr)\in\{0,1\},
\end{equation}
and there exists a time $t_\star$ such that this quantity equals $1$ if and only if $M$ halts on $x$.
Let $\widetilde\rho$ be an imperfect encoding with
$D(\widetilde\rho,\rho_{M,x})\le\delta$.
Then for every $t$,
\begin{equation}
\bigl|\Tr(\Pi_K\,U^t\widetilde\rho\,U^{\dagger t})-\Tr(\Pi_K\,U^t\rho_{M,x}U^{\dagger t})\bigr|\ \le\ \delta,
\end{equation}
and in particular, if $\delta<1/2$ the thresholded predicate
\begin{equation}
\exists t\ge 0:\ \Tr(\Pi_K\,U^t\widetilde\rho\,U^{\dagger t})\ \ge\ 1/2
\end{equation}
is true if and only if $M$ halts on $x$.
\end{proposition}

\begin{proof}
Unitary evolution preserves trace distance:
$D(U^t\widetilde\rho\,U^{\dagger t},U^t\rho_{M,x}U^{\dagger t})=D(\widetilde\rho,\rho_{M,x})$.
Apply Lemma~\ref{lem:proj_lipschitz_trace}.
If $M$ does not halt then the ideal expectation is $0$ for all $t$, hence the imperfect expectation is $<1/2$ for all $t$.
If $M$ halts then at $t=t_\star$ the ideal expectation is $1$, hence the imperfect expectation is $>1/2$.
\end{proof}
